% GENERATED, edit the French unit or the translation table
\mxunit{seconde-s15}{Consolidation and transition to Year 12}\mxsubtitle{25 May--11 June 2027}
\section{Review of calculation, functions and geometry}
\begin{example}[Annual review -- part A]\begin{enumerate}
\item Reduce $168/252$ to lowest terms; calculate $2^{-3}+1/4$.
\item Factorise $x^2-6x+9$, then solve $(x-3)^2=4$.
\item Solve $\dfrac{x-2}{x+1}\geq0$.
\item For $f(x)=4-(x-1)^2$ on $[-2;3]$, give the maximum and the abscissae where $f(x)=0$.
\item With $A(0;1)$, $B(4;3)$, $C(6;4)$, prove collinearity and give an equation of the line.
\end{enumerate}\end{example}
\mxanswerspace[9]{1. $168/252=2/3$, and $2^{-3}+1/4=1/8+2/8=3/8$. 2. $(x-3)^2$; $x-3=\pm2$, so $S=\{1;5\}$. 3. Forbidden $-1$, zero $2$; signs $+,-,+$, so $S=]-\infty;-1[\cup[2;+\infty[$. 4. The square is positive or zero, so $f(x)\leq4$, with equality at $x=1$. The zeros are $x=-1,3$, both in the domain. 5. $\vect{AB}=(4;2)$, $\vect{AC}=(6;3)=\frac32\vect{AB}$. An equation is $x-2y+2=0$.}
\begin{teachernote}Indicative marking: 10 points, 2 per question. For question 3, reserve one point for the domain and the endpoints. For question 5, ask for the conclusion after the calculation.\end{teachernote}

\mxsection{Review of data, probability and modelling}
\begin{example}[Annual review -- part B]\begin{enumerate}
\item Calculate the mean and median of $2;4;4;6;9$. Determine $Q_1,Q_3$.
\item A price drops by $15\,\%$, then rises by $20\,\%$. Give the overall rate and the rate that then restores the initial price.
\item In a fictional population of $100$ people, $40$ play sport ($S$), $30$ come by bike ($V$), $18$ do both. Construct the two-way table and calculate $P_S(V)$, $P_V(S)$ for an equally likely draw.
\item For a rectangle of perimeter $24$ m, express the area as a function of one dimension $x$ and prove its maximum.
\end{enumerate}\end{example}
\mxanswerspace[9]{1. Mean $25/5=5$, median $4$. Quartile ranks $2$ and $4$, so $Q_1=4$, $Q_3=6$.\\2. Overall multiplier $0.85\times1.20=1.02$: an increase of $2\,\%$. Reverse rate $1/1.02-1=-1/51\approx-1.96\,\%$.\\3. Table: $S\cap V=18$, $S\cap\overline V=22$, $\overline S\cap V=12$, $\overline S\cap\overline V=48$. $P_S(V)=18/40=45\,\%$, $P_V(S)=18/30=60\,\%$.\\4. $A(x)=x(12-x)=36-(x-6)^2$ on $]0;12[$. The maximum is $36$ m$^2$, attained for the square of side $6$ m.}
\begin{teachernote}Indicative marking: 10 points, respectively 2, 2, 3, 3. The review of parts A and B can be offered in one hour, then revisited with the course materials.\end{teachernote}

\mxsection{Choosing a strategy and preparing what comes next}
\begin{example}[A program to explain]For a given function $f$, explain the following program. What does it return with $f(x)=x^2$?\end{example}
\begin{lstlisting}
def tableau(f):
    valeurs = []
    for x in range(-3, 4):
        valeurs.append((x, f(x)))
    return valeurs
\end{lstlisting}
\mxanswerspace[3.2]{It returns a list of abscissa--image pairs for the integers from $-3$ to $3$ inclusive: $(-3,9),(-2,4),(-1,1),(0,0),(1,1),(2,4),(3,9)$. The bound $4$ of \texttt{range} is excluded. The table does not describe the images of all real numbers between $-3$ and $3$.}
\begin{keypoint}[Six skills to mobilise]
\textbf{Search}: try things out, organise leads. \textbf{Model}: choose variables, domain and assumptions. \textbf{Represent}: move from text to table, figure or graph. \textbf{Reason}: justify and distinguish conjecture from proof. \textbf{Calculate}: choose a form and check the operations. \textbf{Communicate}: explain clearly with suitable notation and units.
\end{keypoint}
\begin{example}[End-of-year portfolio]Choose a successful piece of work and a corrected one. For each, identify the skills used, explain one method and state a personal goal for the following year. Then redo an exercise matching this goal without looking at the correction.\end{example}
\mxanswerspace[4]{There is no single numerical answer. A satisfactory piece of work precisely identifies the problem, the chosen method, a justification and a check. The goal must be observable, for example: ``state the forbidden values before solving a quotient equation'', rather than ``get better at calculation''.}
\begin{teachernote}To present the Year 12 pathways, rely on the courses actually offered in the school and their programmes in force. The booklet sets no automatic mark threshold for guidance.\end{teachernote}
